% Mean-Field Limit (black-box framework)
% To be included where the main text references ``Mean-Field Limit''.
% Uses \cite{CarmonaDelarueMFGVol1}, \cite{CarmonaDelarueMFGVol2}, \cite{Lacker2018},
% \cite{FournierGuillin2015}, \cite{HuangMalhameCaines2006} --- keys to be added in refs.bib.

\subsection{Mean-Field Limit}\label{sec:mf_limit}

We state a standard common-noise McKean--Vlasov/MFG limit as a black-box framework; our contribution is the constructive solution in the LQG specialization, given in Section~\ref{sec:lqg_constructive}.

\paragraph{Setting.}
Consider a common-noise McKean--Vlasov system on a finite horizon $[0,T]$. Let $\mathcal{F}^0_t := \sigma(W^0_s: s\le t)$ be the filtration generated by a common Brownian motion $W^0$, and for each $N\ge 1$ let $\{B^i\}_{i=1}^N$ be mutually independent and independent of $W^0$. The $N$-agent state processes $\{X^{i,N}_t\}_{i=1}^N$ take values in $\R^d$ and satisfy a coupled system whose coefficients may depend on the empirical measure $\bar\mu^N_t := \frac{1}{N}\sum_{j=1}^N \delta_{X^{j,N}_t}$ and on the common noise. The mean-field limit is described by a McKean--Vlasov SDE in which the conditional law $\mu_t = \Law(X_t\mid \mathcal{F}^0_t)$ of a representative state $X_t$ enters the drift and diffusion; see \cite{CarmonaDelarueMFGVol1,CarmonaDelarueMFGVol2,Lacker2018}.

\paragraph{Assumptions.}
We impose the following standard conditions (see, e.g., \cite{FournierGuillin2015,Lacker2018}).

\begin{enumerate}[label=(H\arabic*),ref=(H\arabic*)]
\item \textbf{(Lipschitz in state and measure.)}
The drift $b(t,x,\mu)$ and diffusion $\sigma(t,x,\mu)$ are Lipschitz in $(x,\mu)$ uniformly in $t\in[0,T]$, with $\mu$ in the space of probability measures on $\R^d$ metrized by the Wasserstein distance $W_2$.

\item \textbf{(Linear growth.)}
$|b(t,x,\mu)| + \|\sigma(t,x,\mu)\| \le C(1+|x| + W_2(\mu,\delta_0))$ for some $C\ge 0$.

\item \textbf{(Common noise integrability.)}
Common-noise coefficients are progressively measurable and square-integrable on $[0,T]$ a.s.

\item \textbf{(Initial condition.)}
The initial states $\{X^{i,N}_0\}_{i=1}^N$ are i.i.d.\ with law $\mu_0$ and satisfy $\E[|X^{i,N}_0|^2]<\infty$; $\mu_0$ has finite second moment.

\item \textbf{(Non-degeneracy / bounded coefficients.)}
The diffusion matrix $\sigma\sigma^\top$ is uniformly elliptic (or at least non-degenerate) and all coefficients are bounded when restricted to compacts in $\R^d\times \mathcal{P}_2(\R^d)$.
\end{enumerate}

Under (H1)--(H5), the limiting McKean--Vlasov equation admits a unique strong solution and the $N$-particle system is well-posed; see \cite{CarmonaDelarueMFGVol1,CarmonaDelarueMFGVol2}.

\begin{theorem}[Convergence of empirical measures]\label{thm:mfg_limit_conv}\label{thm:MFG_fixed_point}
Assume (H1)--(H5). Let $\{X^{i,N}\}_{i=1}^N$ solve the $N$-particle system and let $X$ solve the limiting McKean--Vlasov equation with the same common noise $W^0$ and i.i.d.\ copies $\{B^i\}$ of the idiosyncratic driver. Then for each $t\in[0,T]$,
\[
\bar\mu^N_t \xrightarrow[N\to\infty]{W_2} \mu_t \quad \text{in probability},
\]
where $\mu_t = \Law(X_t\mid \mathcal{F}^0_t)$ is the conditional law of the limit process. In particular, propagation of chaos holds: any fixed finite marginals converge in law to the product of the limit conditional law.
\end{theorem}

\begin{proof}[Proof sketch (black-box references)]
Existence and uniqueness for the limit equation follow from (H1)--(H2) and fixed-point arguments in the space of flows of measures; see \cite{CarmonaDelarueMFGVol1}. Convergence of $\bar\mu^N_t$ to $\mu_t$ in $W_2$ is a standard consequence of (H1)--(H5) and the conditional law of large numbers for the empirical measure; see \cite{FournierGuillin2015} for quantitative rates and \cite{Lacker2018} for the common-noise formulation. Propagation of chaos is then implied by the same estimates; see \cite{CarmonaDelarueMFGVol2}.
\end{proof}

\begin{proposition}[Stability of mean-field equilibrium]\label{prop:mfg_limit_stability}
Under (H1)--(H5), the mean-field equilibrium flow $\{\mu_t\}_{t\in[0,T]}$ is unique in the class of admissible flows with finite second moments. Moreover, $\sup_{t\in[0,T]}\E[W_2(\bar\mu^N_t,\mu_t)^2] \le C/N$ for some constant $C$ depending on $T$ and the data.
\end{proposition}

\begin{proof}[Proof sketch (black-box references)]
Uniqueness follows from the Lipschitz structure (H1) and Gronwall-type estimates in Wasserstein space; see \cite{CarmonaDelarueMFGVol1}. The $O(1/N)$ bound for $\E[W_2(\bar\mu^N_t,\mu_t)^2]$ is given in \cite{FournierGuillin2015} and \cite{HuangMalhameCaines2006} for the classical and MFG settings.
\end{proof}

\begin{remark}[LQG specialization]
In our LQG setting (see Section~IV), the coefficients are affine in state and measure, so (H1)--(H2) hold and the limit is explicit. The preceding black-box statements justify using the mean-field system as the large-$N$ approximation for the particle implementation in Section~VI.
\end{remark}
